
\documentclass{article}
\usepackage{colortbl}
\usepackage{makecell}
\usepackage{multirow}
\usepackage{supertabular}

\begin{document}

\newcounter{utterance}

\centering \large Interaction Transcript for game `cladder', experiment `full\_v1.5\_default', episode 4371 with Bentel rockers\_\_qwen3.5{-}4b\_\_Bentel\_iter\_3.
\vspace{24pt}

{ \footnotesize  \setcounter{utterance}{1}
\setlength{\tabcolsep}{0pt}
\begin{supertabular}{c@{$\;$}|p{.15\linewidth}@{}p{.15\linewidth}p{.15\linewidth}p{.15\linewidth}p{.15\linewidth}p{.15\linewidth}}
    \# & \multicolumn{2}{c}{Player} && \multicolumn{2}{c}{Game Master} \\
    \hline

    \theutterance \stepcounter{utterance}  
    & & & \multicolumn{4}{p{0.6\linewidth}}{
        \cellcolor[rgb]{0.9,0.9,0.9}{
            \makecell[{{p{\linewidth}}}]{
                \texttt{\tiny{[P1$\langle$GM]}}
                \texttt{You are an expert in causal inference. The following question is not a typical commonsense query, but rather a meticulously designed question created by a professor specializing in causal inference, intended to assess the students' mastery of the course content.} \\
\texttt{Imagine a self{-}contained, hypothetical world with only the following conditions, and without any unmentioned factors or causal relationships: Gender has a direct effect on smoking and peanut allergy. Smoking has a direct effect on tar deposit. Tar deposit has a direct effect on peanut allergy. Gender is unobserved. For nonsmokers, the probability of high tar deposit is 28\%. For smokers, the probability of high tar deposit is 58\%. For nonsmokers and with no tar deposit, the probability of being allergic to peanuts is 49\%. For nonsmokers and with high tar deposit, the probability of being allergic to peanuts is 77\%. For smokers and with no tar deposit, the probability of being allergic to peanuts is 27\%. For smokers and with high tar deposit, the probability of being allergic to peanuts is 62\%. The overall probability of smoking is 58\%. Will smoking increase the chance of being allergic to peanuts?} \\
\texttt{Start your answer with 'Yes' or 'No', followed by additional reasoning or evidence to support your explanation.} \\
            }
        }
    }
     \\ \\

    \theutterance \stepcounter{utterance}  
    & \multicolumn{4}{p{0.6\linewidth}}{
        \cellcolor[rgb]{0.561,0.737,0.561}{
            \makecell[{{p{\linewidth}}}]{
                \texttt{\tiny{[P1$\rangle$GM]}}
                \texttt{Yes} \\
\\ 
\texttt{To determine whether smoking increases the chance of being allergic to peanuts, we must calculate the counterfactual probability: the probability of peanut allergy for a person *if they had not smoked*, while holding all other causal factors constant (specifically, the distribution of tar deposits and the underlying gender distribution). In causal inference, this requires adjusting for confounding variables. Here, **Gender** is an unobserved confounder that directly affects both **Smoking** and **Peanut Allergy**, and **Tar Deposit** is a mediator through which smoking influences allergy.} \\
\\ 
\texttt{Because Gender is unobserved, we cannot simply compare the observed prevalence of smoking and allergy. Instead, we must compute the potential outcome \$P(\textbackslash{}text\{Allergy\} | \textbackslash{}text\{No Smoking\})\$ by integrating over the unobserved gender distribution and the induced distribution of tar deposits in the absence of smoking.} \\
\\ 
\texttt{Let \$G\$ denote Gender, \$S\$ denote Smoking, \$T\$ denote Tar Deposit, and \$A\$ denote Allergy.} \\
\texttt{We are given:} \\
\texttt{1.  \$P(S) = 0.58\$, so \$P(\textbackslash{}neg S) = 0.42\$.} \\
\texttt{2.  \$P(T|S, \textbackslash{}neg G) = 0.58\$, \$P(T|S, G) = 0.58\$ (The prompt implies the smoking{-}tar relationship is uniform across genders, or at least the marginal \$P(T|S)\$ is given as 0.58. Let's assume the conditional probabilities for tar given smoking are the same regardless of gender, or that the "58\%" applies to the aggregate smoker population. Similarly for non{-}smokers, \$P(T|\textbackslash{}neg S) = 0.28\$).} \\
\texttt{3.  \$P(A|\textbackslash{}neg S, \textbackslash{}neg T, \textbackslash{}neg G) = 0.49\$, \$P(A|\textbackslash{}neg S, T, \textbackslash{}neg G) = 0.77\$.} \\
\texttt{4.  \$P(A|S, \textbackslash{}neg T, \textbackslash{}neg G) = 0.27\$, \$P(A|S, T, \textbackslash{}neg G) = 0.62\$.} \\
\texttt{5.  \$P(A|\textbackslash{}neg S, \textbackslash{}neg T, G) = 0.49\$, \$P(A|\textbackslash{}neg S, T, G) = 0.77\$.} \\
\texttt{6.  \$P(A|S, \textbackslash{}neg T, G) = 0.27\$, \$P(A|S, T, G) = 0.62\$.} \\
\texttt{    *(Note: The prompt provides probabilities for "nonsmokers" and "smokers" generally, but the specific values for "nonsmokers and with no tar" etc., are identical for both genders in the prompt's phrasing, implying the effect of Tar on Allergy is independent of Gender, or the provided numbers are the marginal probabilities within those smoking/non{-}smoking groups. Let's assume the effect of Tar on Allergy is constant across genders based on the identical numbers provided).*} \\
\\ 
\texttt{To find the causal effect of smoking, we need to compare the observed \$P(A|S)\$ with the counterfactual \$P(A|\textbackslash{}neg S)\$. However, since Gender confounds the relationship between Smoking and Allergy (it affects both directly and indirectly via Tar), the observed \$P(A|S)\$ is a weighted average of the effects in men and women. The counterfactual \$P(A|\textbackslash{}neg S)\$ will also be a weighted average, but the weights (the proportion of men and women among non{-}smokers) might differ from the weights among smokers if Gender affects the propensity to smoke.} \\
\\ 
\texttt{Let \$p = P(G)\$ be the proportion of the population that is Gender A (e.g., Male), and \$1{-}p\$ be Gender B (e.g., Female).} \\
\texttt{Let \$s\_m = P(S|G)\$ and \$s\_f = P(S|\textbackslash{}neg G)\$ be the propensity to smoke for each gender.} \\
\texttt{We know \$P(S) = p \textbackslash{}cdot s\_m + (1{-}p) \textbackslash{}cdot s\_f = 0.58\$.} \\
\\ 
\texttt{The probability of Tar given Smoking is given as 0.58 for the aggregate. Assuming the effect of smoking on tar is identical for both genders (as no gender{-}specific tar probabilities are listed), then \$P(T|S, G) = 0.58\$ and \$P(T|S, \textbackslash{}neg G) = 0.58\$.} \\
\texttt{Similarly, \$P(T|\textbackslash{}neg S, G) = 0.28\$ and \$P(T|\textbackslash{}neg S, \textbackslash{}neg G) = 0.28\$.} \\
\\ 
\texttt{The probability of Allergy given Smoking (\$P(A|S)\$) is:} \\
\texttt{\$P(A|S) = P(A|S, T)P(T|S) + P(A|S, \textbackslash{}neg T)P(\textbackslash{}neg T|S)\$} \\
\texttt{\$P(A|S) = 0.62(0.58) + 0.27(0.42) = 0.3596 + 0.1134 = 0.473\$.} \\
\\ 
\texttt{Now, we calculate the counterfactual \$P(A|\textbackslash{}neg S)\$. This is the probability of allergy if everyone had not smoked.} \\
\texttt{\$P(A|\textbackslash{}neg S) = P(A|\textbackslash{}neg S, T)P(T|\textbackslash{}neg S) + P(A|\textbackslash{}neg S, \textbackslash{}neg T)P(\textbackslash{}neg T|\textbackslash{}neg S)\$} \\
\texttt{\$P(A|\textbackslash{}neg S) = 0.77(0.28) + 0.49(0.72) = 0.2156 + 0.3528 = 0.5684\$.} \\
\\ 
\texttt{Wait, this calculation assumes that the distribution of Tar in the non{-}smoking population is simply determined by the rule "Nonsmokers have 28\% high tar," regardless of gender. The prompt states: "For nonsmokers, the probability of high tar deposit is 28\%." This is a marginal statement. However, in a causal world with an unobserved confounder (Gender), the marginal probability \$P(T|\textbackslash{}neg S)\$ depends on the distribution of Gender among non{-}smokers.} \\
\\ 
\texttt{Let's re{-}evaluate carefully.} \\
\texttt{Let \$G \textbackslash{}in \textbackslash{}\{0, 1\textbackslash{}\}\$.} \\
\texttt{\$P(S|G=0) = s\_0\$, \$P(S|G=1) = s\_1\$.} \\
\texttt{\$P(S) = p s\_0 + (1{-}p) s\_1 = 0.58\$.} \\
\texttt{\$P(T|S, G) = 0.58\$ (Assuming uniform effect of smoking on tar).} \\
\texttt{\$P(T|\textbackslash{}neg S, G) = 0.28\$ (Assuming uniform effect of non{-}smoking on tar).} \\
\\ 
\texttt{If the effect of smoking on tar is truly uniform across genders, then the distribution of Tar given Smoking is 0.58 regardless of Gender.} \\
\texttt{Consequently, the distribution of Tar given Non{-}Smoking is 0.28 regardless of Gender.} \\
\texttt{In this specific hypothetical setup, the unobserved confounder (Gender) does not create a *difference* in the conditional probabilities of the mediator (Tar) or the outcome (Allergy) between the two genders, *given* the smoking status. The prompt gives specific probabilities for "nonsmokers" and "smokers" without distinguishing gender, implying these are the conditional probabilities \$P(T|S)\$ and \$P(A|S, T)\$ that hold universally or that the gender{-}specific parameters are identical.} \\
\\ 
\texttt{If the parameters for Tar and Allergy are identical for both genders, then the counterfactual \$P(A|\textbackslash{}neg S)\$ is simply the calculation performed above:} \\
\texttt{\$P(A|\textbackslash{}neg S) = P(A|\textbackslash{}neg S, T)P(T|\textbackslash{}neg S) + P(A|\textbackslash{}neg S, \textbackslash{}neg T)P(\textbackslash{}neg T|\textbackslash{}neg S)\$} \\
\texttt{\$P(A|\textbackslash{}neg S) = 0.77 \textbackslash{}times 0.28 + 0.49 \textbackslash{}times 0.72\$} \\
\texttt{\$P(A|\textbackslash{}neg S) = 0.2156 + 0.3528 = 0.5684\$.} \\
\\ 
\texttt{Comparing the observed probability with smoking (\$0.473\$) to the counterfactual probability without smoking (\$0.5684\$):} \\
\texttt{\$0.473 < 0.5684\$.} \\
\\ 
\texttt{Therefore, smoking decreases the chance of being allergic to peanuts in this specific hypothetical world. The mechanism is that smoking increases the probability of high tar deposits (from 28\% to 58\%), and high tar deposits *increase* the probability of peanut allergy (from 49\% to 77\% for non{-}smokers, and from 27\% to 62\% for smokers). Since the path \$S \textbackslash{}to T \textbackslash{}to A\$ is positive (Smoking \$\textbackslash{}uparrow\$ Tar \$\textbackslash{}uparrow\$ Allergy), and there are no negative direct effects of Smoking on Allergy mentioned (the prompt only gives probabilities conditional on T, implying the direct effect is zero or subsumed in the T{-}conditioned probabilities), the net effect of smoking is to increase the likelihood of the intermediate state (High Tar), which in turn increases the likelihood of the outcome (Allergy).} \\
\\ 
\texttt{Wait, let me double check the "direct effect" wording. "Gender has a direct effect on smoking and peanut allergy." "Smoking has a direct effect on tar deposit." "Tar deposit has a direct effect on peanut allergy."} \\
\texttt{This implies the causal graph:} \\
\texttt{\$G \textbackslash{}to S\$} \\
\texttt{\$G \textbackslash{}to A\$} \\
\texttt{\$S \textbackslash{}to T\$} \\
\texttt{\$T \textbackslash{}to A\$} \\
\\ 
\texttt{The question asks: "Will smoking increase the chance of being allergic to peanuts?"} \\
\texttt{This is asking for the sign of the causal effect \$P(A|do(S=1)) {-} P(A|do(S=0))\$.} \\
\texttt{Using the backdoor adjustment or simply the causal path decomposition:} \\
\texttt{The total effect of \$S\$ on \$A\$ is the sum of the direct effect (if any, not explicitly quantified as separate from the T{-}path) and the indirect effect through \$T\$.} \\
\texttt{The prompt gives us \$P(A|S, T)\$ and \$P(T|S)\$. It does not explicitly give a "direct effect of S on A" separate from the T{-}path. Usually, in such problems, the conditional probability \$P(A|S, T)\$ represents the residual risk after accounting for T. If there were a direct effect \$S \textbackslash{}to A\$ independent of \$T\$, it would likely be specified or the numbers would be inconsistent with a pure mediation model.} \\
\texttt{Let's assume the standard mediation model where the effect of S on A is entirely mediated by T, or that the provided numbers encapsulate the total conditional probability.} \\
\texttt{Effect = \$E{[}A|S=1{]} {-} E{[}A|S=0{]}\$.} \\
\texttt{We calculated \$E{[}A|S=1{]} = 0.473\$.} \\
\texttt{We calculated \$E{[}A|S=0{]} = 0.5684\$.} \\
\texttt{The difference is \$0.473 {-} 0.5684 = {-}0.0954\$.} \\
\texttt{The effect is negative. Smoking *decreases* the chance of peanut allergy in this world.} \\
\\ 
\texttt{Why? Because the "bad" thing (High Tar) is more likely in smokers (58\%) than non{-}smokers (28\%), and High Tar is a "bad" thing for allergies (increases probability from \textasciitilde{}50\% to \textasciitilde{}70\%). Thus, by causing High Tar, Smoking causes Allergies.} \\
\texttt{Wait, my logic in the previous step was: "Smoking increases the probability of high tar deposits... and high tar deposits increase the probability of peanut allergy."} \\
\texttt{If High Tar increases Allergy, and Smoking increases High Tar, then Smoking should *increase* Allergy.} \\
\texttt{Let's re{-}calculate the numbers.} \\
\texttt{Non{-}smoker, No Tar: 49\%} \\
\texttt{Non{-}smoker, High Tar: 77\%} \\
\texttt{Smoker, No Tar: 27\%} \\
\texttt{Smoker, High Tar: 62\%} \\
\\ 
\texttt{Base rate of High Tar in Non{-}smokers: 28\%.} \\
\texttt{Base rate of High Tar in Smokers: 58\%.} \\
\\ 
\texttt{\$P(A|\textbackslash{}neg S) = 0.49(1{-}0.28) + 0.77(0.28) = 0.49(0.72) + 0.77(0.28) = 0.3528 + 0.2156 = 0.5684\$.} \\
\texttt{\$P(A|S) = 0.27(1{-}0.58) + 0.62(0.58) = 0.27(0.42) + 0.62(0.58) = 0.1134 + 0.3596 = 0.4730\$.} \\
\\ 
\texttt{Result: \$P(A|S) < P(A|\textbackslash{}neg S)\$.} \\
\texttt{So, despite the mechanism "Smoking {-}> High Tar {-}> Allergy" seeming positive, the *conditional probabilities* provided show that for smokers, having NO tar deposit is associated with a much lower allergy rate (27\%) compared to non{-}smokers with no tar (49\%).} \\
\texttt{This implies that the "No Tar" state is much more protective for smokers than for non{-}smokers.} \\
\texttt{Why? Perhaps because for non{-}smokers, the baseline risk of allergy (49\%) is high even without tar, whereas for smokers, the baseline risk (27\%) is low.} \\
\texttt{The "direct effect" of Gender might be the key here. Gender causes Smoking and Allergy.} \\
\texttt{If Gender = Male: High propensity to smoke? High propensity to allergy?} \\
\texttt{If Gender = Female: Low propensity to smoke? Low propensity to allergy?} \\
\texttt{If the "No Tar" probability of allergy is 49\% for non{-}smokers generally, but 27\% for smokers generally, this suggests that the group of non{-}smokers has a higher inherent risk of allergy than the group of smokers. Since Gender causes both, the "No Tar" group of non{-}smokers is likely composed of a gender with high baseline allergy risk, while the "No Tar" group of smokers is composed of a gender with low baseline allergy risk.} \\
\texttt{However, the question is simply "Will smoking increase the chance...?" based on the provided numbers.} \\
\texttt{The calculation shows \$P(A|S) = 0.473\$ and \$P(A|\textbackslash{}neg S) = 0.568\$.} \\
\texttt{Since \$0.473 < 0.568\$, the answer is **No**. Smoking decreases the chance.} \\
\\ 
\texttt{Let's double check the arithmetic.} \\
\texttt{\$0.49 \textbackslash{}times 0.72 = 0.3528\$} \\
\texttt{\$0.77 \textbackslash{}times 0.28 = 0.2156\$} \\
\texttt{Sum = \$0.5684\$. Correct.} \\
\\ 
\texttt{\$0.27 \textbackslash{}times 0.42 = 0.1134\$} \\
\texttt{\$0.62 \textbackslash{}times 0.58 = 0.3596\$} \\
\texttt{Sum = \$0.4730\$. Correct.} \\
\\ 
\texttt{The difference is \$0.473 {-} 0.5684 = {-}0.0954\$.} \\
\texttt{The causal effect is negative.} \\
\texttt{Therefore, smoking will **not** increase the chance; it will decrease it.} \\
\\ 
\texttt{The confusion in my initial intuitive check ("Smoking increases Tar, Tar increases Allergy, so Smoking increases Allergy") was ignoring the massive difference in the baseline probabilities (\$P(A|\textbackslash{}neg T)\$) between smokers and non{-}smokers.} \\
\texttt{For non{-}smokers: \$P(A|\textbackslash{}neg T) = 49\textbackslash{}\%\$.} \\
\texttt{For smokers: \$P(A|\textbackslash{}neg T) = 27\textbackslash{}\%\$.} \\
\texttt{Even though smoking pushes you into the "High Tar" bucket (which raises risk to 62\% vs 77\% for non{-}smokers), the "No Tar" bucket for smokers is so much safer (27\% vs 49\%) that the overall average drops.} \\
\texttt{The "direct effect" of Gender is likely driving this: The gender that smokes less (and thus falls into the non{-}smoker category) has a much higher inherent risk of peanut allergy than the gender that smokes more.} \\
\texttt{Regardless of the underlying gender distribution, the counterfactual \$P(A|\textbackslash{}neg S)\$ is fixed by the rules of the world: "For nonsmokers, the probability of high tar is 28\%" and the subsequent allergy probabilities. The counterfactual world where everyone is a non{-}smoker has a 28\% tar rate and the associated allergy rates. The observed world where 58\% smoke has a 58\% tar rate and the associated allergy rates.} \\
\texttt{The comparison is strictly between the aggregate of the non{-}smoking world and the aggregate of the smoking world.} \\
\texttt{\$P(A|\textbackslash{}neg S) = 0.5684\$.} \\
\texttt{\$P(A|S) = 0.4730\$.} \\
\texttt{Smoking reduces the probability.} \\
\\ 
\texttt{Answer: No.} \\
\\ 
\texttt{Reasoning Summary:} \\
\texttt{1.  **Define the Counterfactual**: To determine if smoking increases the chance of peanut allergy, we compare the observed probability of allergy given smoking (\$P(A|S)\$) with the counterfactual probability of allergy given no smoking (\$P(A|\textbackslash{}neg S)\$).} \\
\texttt{2.  **Calculate \$P(A|\textbackslash{}neg S)\$**: In a world where no one smokes, the probability of high tar deposit is 28\%.} \\
\texttt{    *   \$P(A|\textbackslash{}neg S, \textbackslash{}text\{No Tar\}) = 49\textbackslash{}\%\$} \\
\texttt{    *   \$P(A|\textbackslash{}neg S, \textbackslash{}text\{High Tar\}) = 77\textbackslash{}\%\$} \\
\texttt{    *   \$P(A|\textbackslash{}neg S) = 0.49 \textbackslash{}times (1 {-} 0.28) + 0.77 \textbackslash{}times 0.28 = 0.49 \textbackslash{}times 0.72 + 0.77 \textbackslash{}times 0.28 = 0.3528 + 0.2156 = 0.5684\$ (or 56.84\%).} \\
\texttt{3.  **Calculate \$P(A|S)\$**: In the observed world where 58\% smoke, the probability of high tar deposit is 58\%.} \\
\texttt{    *   \$P(A|S, \textbackslash{}text\{No Tar\}) = 27\textbackslash{}\%\$} \\
\texttt{    *   \$P(A|S, \textbackslash{}text\{High Tar\}) = 62\textbackslash{}\%\$} \\
\texttt{    *   \$P(A|S) = 0.27 \textbackslash{}times (1 {-} 0.58) + 0.62 \textbackslash{}times 0.58 = 0.27 \textbackslash{}times 0.42 + 0.62 \textbackslash{}times 0.58 = 0.1134 + 0.3596 = 0.4730\$ (or 47.30\%).} \\
\texttt{4.  **Compare**: The probability of peanut allergy is lower in the smoking group (47.30\%) than in the non{-}smoking group (56.84\%).} \\
\texttt{5.  **Conclusion**: Although smoking increases the likelihood of high tar deposits (which are associated with higher allergy rates), the baseline risk of allergy for nonsmokers (49\%) is significantly higher than for smokers (27\%). This suggests a strong confounding effect by Gender (where the gender less likely to smoke has a higher inherent allergy risk). Consequently, the net causal effect of smoking is to *decrease* the probability of peanut allergy.} \\
\\ 
\texttt{Final Answer Formulation:} \\
\texttt{Start with "No".} \\
\texttt{Explain the calculation of the counterfactual probability.} \\
\texttt{Show the comparison.} \\
\texttt{Conclude that the effect is negative.} \\
\\ 
\texttt{One detail: "Gender is unobserved." Does this change the counterfactual calculation?} \\
\texttt{The counterfactual \$P(A|\textbackslash{}neg S)\$ is the probability of A in a world where \$S\$ is set to 0 for everyone.} \\
\texttt{In this new world, the distribution of Gender remains the same as the original world (causal interventions on S do not change the distribution of G).} \\
\texttt{The rules for the world are:} \\
\texttt{{-} If \$S=0\$, \$P(T=1) = 0.28\$.} \\
\texttt{{-} If \$S=0, T=0\$, \$P(A=1) = 0.49\$ (for both genders, as the numbers are identical).} \\
\texttt{{-} If \$S=0, T=1\$, \$P(A=1) = 0.77\$ (for both genders).} \\
\texttt{Therefore, regardless of the unknown gender distribution \$p\$, the aggregate \$P(A|\textbackslash{}neg S)\$ is simply the weighted average of the two conditional probabilities using the fixed \$P(T|\textbackslash{}neg S) = 0.28\$. The gender distribution cancels out because the conditional probabilities \$P(A|S, T, G)\$ and \$P(T|S, G)\$ are identical for both genders in the prompt's data.} \\
\texttt{Thus, the calculation holds.} \\
\\ 
\texttt{The answer is definitively No.} \\
\\ 
\texttt{Structure of the response:} \\
\texttt{1.  **No**} \\
\texttt{2.  **Reasoning**:} \\
\texttt{    *   Define the goal: Compare \$P(A|do(S=1))\$ vs \$P(A|do(S=0))\$.} \\
\texttt{    *   Calculate \$P(A|S=0)\$: Use the given non{-}smoker stats. \$0.49 \textbackslash{}times 0.72 + 0.77 \textbackslash{}times 0.28 = 0.5684\$.} \\
\texttt{    *   Calculate \$P(A|S=1)\$: Use the given smoker stats. \$0.27 \textbackslash{}times 0.42 + 0.62 \textbackslash{}times 0.58 = 0.4730\$.} \\
\texttt{    *   Interpret the result: Since \$0.4730 < 0.5684\$, smoking reduces the probability.} \\
\texttt{    *   Explain the intuition: The "direct effect" of the unobserved confounder (Gender) creates a situation where the non{-}smoking group has a much higher baseline risk of allergy than the smoking group, outweighing the increased risk from tar deposits.} \\
\\ 
\texttt{This seems complete and correct.} \\
\\ 
\texttt{One minor check: "For nonsmokers and with no tar deposit, the probability of being allergic to peanuts is 49\%."} \\
\texttt{Does this apply to both genders? The prompt says "For nonsmokers...". It does not say "For nonsmokers of Gender X...". Given the symmetry of the other numbers (Smoker/No Tar 27\% vs Nonsmoker/No Tar 49\%), and the lack of gender{-}specific sub{-}parameters, it is standard interpretation in such problems that the conditional probabilities provided are the aggregate or uniform across the unobserved confounder. If they varied by gender, the problem would be unsolvable without knowing the gender distribution. Since it asks for a definitive "Yes/No", we must assume the parameters are uniform or the aggregation works out such that the counterfactual is well{-}defined by the marginal statements. The uniformity assumption is the only} \\
            }
        }
    }
    & & \\ \\

    \theutterance \stepcounter{utterance}  
    & & & \multicolumn{2}{p{0.3\linewidth}}{
        \cellcolor[rgb]{0.9,0.9,0.9}{
            \makecell[{{p{\linewidth}}}]{
                \texttt{\tiny{[GM$|$GM]}}
                \texttt{Yes} \\
            }
        }
    }
    & & \\ \\

    \theutterance \stepcounter{utterance}  
    & & & \multicolumn{2}{p{0.3\linewidth}}{
        \cellcolor[rgb]{0.9,0.9,0.9}{
            \makecell[{{p{\linewidth}}}]{
                \texttt{\tiny{[GM$|$GM]}}
                \texttt{yes} \\
            }
        }
    }
    & & \\ \\

    \theutterance \stepcounter{utterance}  
    & & & \multicolumn{2}{p{0.3\linewidth}}{
        \cellcolor[rgb]{0.9,0.9,0.9}{
            \makecell[{{p{\linewidth}}}]{
                \texttt{\tiny{[GM$|$GM]}}
                \texttt{game\_result = WIN} \\
            }
        }
    }
    & & \\ \\

\end{supertabular}
}

\end{document}
